\documentclass[11pt,a4paper]{article}
\input{../../../sablony/pripravy/priprava.sty}
\pagestyle{empty}
\geometry{a4paper,top=16mm,bottom=17mm,left=22mm,right=22mm}
\begin{document}
{\Large\bfseries Pracovní list 71 – výsledky pro učitele}\par
\smallskip
\textbf{1.} $f(x)=|x-4|+|x+4|$. Body změny $-4$ a $4$.
\[
\begin{array}{c|ccc}
 &x<-4&-4\le x\le4&x>4\\\hline
|x-4|&4-x&4-x&x-4\\
|x+4|&-x-4&x+4&x+4\\\hline
f(x)&-2x&8&2x
\end{array}
\]
\[
f(x)=\begin{cases}-2x,&x<-4,\\8,&-4\le x\le4,\\2x,&x>4.\end{cases}
\]
$D(f)=\mathbb R$, $H(f)=\langle8;+\infty)$, minimum $8$ nastává na $\langle-4;4\rangle$. Graf: levá polopřímka přes $[-4;8]$, vodorovná úsečka do $[4;8]$, pravá polopřímka.
\bigskip
\textbf{2.} $p(x)=|x+2|+|x-2|$:
\[
p(x)=\begin{cases}-2x,&x<-2,\\4,&-2\le x\le2,\\2x,&x>2.\end{cases}
\]
$D(p)=\mathbb R$, $H(p)=\langle4;+\infty)$, interval konstantnosti $\langle-2;2\rangle$, minimum $4$.
\bigskip
\textbf{3.} $q(x)=|x-2|-|x+1|$:
\[
q(x)=\begin{cases}3,&x<-1,\\1-2x,&-1\le x\le2,\\-3,&x>2.\end{cases}
\]
$D(q)=\mathbb R$, $H(q)=\langle-3;3\rangle$. Krajní body prostřední úsečky: $[-1;3]$, $[2;-3]$.
\bigskip
\textbf{4.} $r(x)=|2x-2|+|x+2|$: změny při $x=-2$ a $x=1$.
\[
r(x)=\begin{cases}-3x,&x<-2,\\4-x,&-2\le x\le1,\\3x,&x>1.\end{cases}
\]
Minimum $3$ v $x=1$. Střední část není konstantní; první člen má dvojnásobnou strmost oproti $|x-1|$.
\medskip
\textbf{Typické chyby:} změna znaménka před celou závorkou (např. $-(x-4)=4-x$); chybné pořadí nulových bodů; chybějící dosazení $x=\pm4$; zaměnění celé přímky za její platnou část.
\end{document}
